\chapter{Stakeholder und Use Cases}

\begin{lernziele}
Nach diesem Kapitel können Sie Stakeholder identifizieren und aus typischen Nutzungsszenarien Anforderungen ableiten.
\end{lernziele}

\section{Stakeholder identifizieren}
Stakeholder sind nicht nur Endnutzer. Auch Betreiber, Wartung, Entwicklung, Management und Sicherheitsverantwortliche haben Interessen. Ein Roboter, der technisch funktioniert, aber schlecht wartbar ist, kann im Betrieb trotzdem scheitern.

\section{Stakeholder des Roboters}
\begin{itemize}
\item Nutzer: möchte Ziele einfach vorgeben.
\item Betreiber: möchte zuverlässigen Betrieb.
\item Wartungstechniker: benötigt Diagnoseinformationen.
\item Entwickler: benötigt klare Schnittstellen.
\item Sicherheitsverantwortlicher: achtet auf Risiken.
\end{itemize}

\section{Use Cases}
Use Cases beschreiben, wie ein Akteur mit dem System interagiert. Sie sind keine detaillierten Algorithmen, sondern Nutzungsbeschreibungen.

\section{Beispiel: Ziel anfahren}
\textbf{Akteur:} Nutzer\\
\textbf{Ziel:} Roboter fährt zu einem ausgewählten Zielpunkt.\\
\textbf{Vorbedingung:} Roboter ist eingeschaltet und lokalisiert.\\
\textbf{Ablauf:}
\begin{enumerate}
\item Nutzer wählt Ziel.
\item Roboter prüft Verfügbarkeit der Karte.
\item Roboter plant Pfad.
\item Roboter fährt los.
\item Roboter weicht Hindernissen aus.
\item Roboter erreicht Ziel.
\end{enumerate}

\section{Use Cases und Anforderungen}
Aus einem Use Case lassen sich Anforderungen ableiten. Wenn der Roboter Ziele anfahren soll, braucht er Navigation, Lokalisierung, Pfadplanung und Hindernisvermeidung.

\begin{uebungbox}
Schreiben Sie einen Use Case \enquote{Zur Ladestation zurückkehren}. Leiten Sie daraus mindestens drei Anforderungen ab.
\end{uebungbox}
